\section{SECCIÓN 4. DOCUMENTACIÓN}
La documentación que describe y respalda el software del proyecto o el proceso de desarrollo de software se creará y actualizará periódicamente durante todo el ciclo de desarrollo.

La Tabla 4-1 es una lista de los productos de software entregables del proyecto y el estándar o pautas asociados utilizados para desarrollar y mantener los productos de software. Cualquier pauta de adaptación también se encuentra en la Tabla 4-1. La Tabla 4-2 es una lista de productos no entregables.

Para los documentos de software del proyecto que se desarrollarán y que aún no figuran en las Tablas 4-1 y 4-2, SQA ayudará a identificar las especificaciones, estándares y Descripciones de Elementos de Datos (DIDs) que se seguirán en la preparación de la documentación requerida.

\begin{table}[H]
\centering
\caption{Productos de Software Entregables}
\label{tab:deliverable}
\begin{tabular}{|p{3cm}|p{4cm}|p{6cm}|}
\hline
\textbf{NOMENCLATURA} & \textbf{DOCUMENTACIÓN ENTREGABLE} & \textbf{DID, ESTÁNDAR, PAUTA} \\
\hline
[Nombre del CI] & [TIPO DE DOCUMENTO, ej., SSS] & [DID, ej., DI-IPSC-81431 de MIL-STD-498] \\
\hline
[Nombre del CI] & [TIPO DE DOCUMENTO] & [DID, ESTÁNDAR, PAUTA] \\
\hline
[Nombre del CI] & [TIPO DE DOCUMENTO] & [DID, ESTÁNDAR, PAUTA] \\
\hline
\end{tabular}
\end{table}

\begin{table}[H]
\centering
\caption{Productos de Software No Entregables}
\label{tab:nondeliverable}
\begin{tabular}{|c|}
\hline
\textbf{TÍTULO DEL DOCUMENTO} \\
\hline
[Título del documento] \\
\hline
[Título del documento] \\
\hline
[Título del documento] \\
\hline
\end{tabular}
\end{table}

\textit{Indicar cómo se verificará la adecuación de los documentos. El proceso de revisión de documentos debe incluir los criterios y la identificación de la revisión o auditoría mediante la cual se confirmará la adecuación de cada documento.}

\begin{enumerate}
    \item Todos los documentos se someterán a una revisión por pares de acuerdo con el Proceso de Revisión por Pares, referencia (n).
\end{enumerate}

Tras la finalización de una revisión por pares, SQA registra y reporta las mediciones de la revisión por pares (el elemento revisado, el número de errores detectados, la fase en la que se realizó la revisión por pares, el número de informes de error cerrados y el número de informes de error abiertos) a SEPO.

Tras la finalización de una revisión por pares, el producto de software se enviará a SCM y se colocará bajo control de CM. El producto de software se procesará entonces de acuerdo con el proceso de aprobación y liberación de productos de software de SCM descrito en la referencia (e).
